\chapter{Requisiti}

Il lavoro di semestre si inserisce nel contesto del paper \emph{Branch-and-Bound Algorithms as Polynomial-time Approximation Scheme}~\cite{bb-ptas}, scritto dal nostro relatore Prof.\ Monaldo Mastrolilli insieme a István Encz Koppányi ed Eleonora Vercesi. Gli autori forniscono un’implementazione Python del Branch-and-Bound per il caso delle \emph{unrelated machines}, comprensiva di euristiche e gestione dei bound. Il nostro compito consiste nel comprendere a fondo tale implementazione, adattarla al caso delle \emph{identical machines} e integrare la tecnica di \emph{profiling}, applicabile poiché le identical machines rappresentano un caso particolare di uniform machines. L’obiettivo finale è verificare sperimentalmente quanto presentato nel paper, confrontando la versione base dell’algoritmo con quella dotata di profiling.

\section{Requisiti funzionali}

Il progetto richiede i seguenti requisiti funzionali:

\begin{description}

    \item[RF1 -- Solver per identical machines]~\\ Realizzare una variante dell’algoritmo Branch-and-Bound in grado di risolvere il problema di scheduling su identical machines, adattando la struttura e le funzioni dell’implementazione esistente per le unrelated machines.

    \item[RF2 -- Integrazione del profiling]~\\ Implementare il meccanismo di profiling descritto nel paper, così da ottenere una seconda variante dell’algoritmo direttamente confrontabile con la versione base.

    \item[RF3 -- Generazione ed esecuzione degli esperimenti]~\\ Eseguire entrambe le varianti dell’algoritmo e registrare i risultati, raccogliendo una quantità consona di dati.

    \item[RF4 -- Raccolta dei dati]~\\ Memorizzare in formato strutturato informazioni quali nodi esplorati, tempo di esecuzione, delta rispetto alla soluzione ottima dell’istanza considerata e altre metriche rilevanti.

    \item[RF5 -- Validazione sperimentale]~\\ Analizzare i dati raccolti, con l’aiuto di strumenti statistici e di visualizzazione, per verificare l’efficacia del profiling rispetto alla versione base dell’algoritmo.

\end{description}

\section{Requisiti non funzionali}

Non erano previsti requisiti non funzionali formali, ma il progetto implicava alcuni vincoli pratici:

\begin{description}

    \item[Efficienza computazionale]~\\ L’implementazione e la scelta delle istanze dovevano permettere di generare e analizzare un numero significativo di dati entro i limiti temporali del lavoro di semestre.

    \item[Ripetibilità]~\\ La generazione delle istanze e l’esecuzione degli algoritmi dovevano essere riproducibili tramite seed controllati.

    \item[Mantenibilità e pulizia del codice]~\\ Anche trattandosi di un esperimento non destinato a evoluzione futura, il codice doveva essere scritto in modo chiaro e integrarsi coerentemente con la struttura esistente.

\end{description}
